% Part: many-valued-logic
% Chapter: syntax-and-semantics
% Section: matrices

\documentclass[../../../include/open-logic-section]{subfiles}

\begin{document}

\olfileid{mvl}{syn}{mat}

\olsection{మాత్రికలు}

ఒక బహుమూల్య తర్కాన్ని దాని భాష, సత్యమూల్యాల సమితి~$V$,
నిర్దేశిత సత్యమూల్యాల ఉపసమితి, దాని సంయోజకాలకు చెందిన
సత్యమూల్య ప్రమేయాలు నిర్ణయిస్తాయి. ఈ అంశాలన్నిటినీ
కలిపి ఒక \emph{మాత్రిక} అంటాం.

\begin{defn}[మాత్రిక]
\ollabel{defn:matrix}
తర్కం~$\Log L$కు ఒక \emph{మాత్రిక}లో కింది అంశాలు ఉంటాయి:
\begin{enumerate}
\item భాష~$\Lang L$ను ఏర్పరచే సంయోజకాల సమితి;
\item ఖాళీ కాని సత్యమూల్యాల సమితి $V \neq \emptyset$;
\item నిర్దేశిత సత్యమూల్యాల సమితి $V^+ \subseteq V$;
\item $\Lang L$లోని ప్రతి $n$-స్థానిక సంయోజకం~$\star$కు
సత్యమూల్య ప్రమేయం~$\tf{\star} : V^n \to V$.
$n = 0$ అయితే, $\tf{\star}$ అనేది $V$లోని ఒక మూలకం మాత్రమే.
\end{enumerate}
\end{defn}

\begin{ex}
సాంప్రదాయిక తర్కం~\LogCL{} యొక్క మాత్రికలో కింది అంశాలు ఉంటాయి:
\begin{enumerate}
  \item $\lfalse$, $\lnot$, $\land$, $\lor$, $\lif$ గల
  ప్రామాణిక ప్రతిజ్ఞావాక్య భాష~$\Lang L_0$.
  \item సత్యమూల్యాల సమితి $V = \{\True, \False\}$.
  \item $\True$ ఒక్కటే నిర్దేశిత విలువ; అంటే $V^+ = \{\True\}$.
  \item $\lfalse$కు $\tf{\lfalse} = \False$. ఇతర
  సత్యమూల్య ప్రమేయాలను సాధారణ సత్య పట్టికలు ఇస్తాయి
  (\olref{fig:tf-CL} చూడండి).
\end{enumerate}
\begin{figure}
  \begin{center}
      \begin{tabular}{c|c}
        $\tf{\lnot}$ & \\
        \hline
        $\True$ & $\False$ \\
        $\False$ & $\True$
      \end{tabular}
      \quad
      \begin{tabular}{c|cc}
        $\tf{\land}$ & $\True$ & $\False$ \\
        \hline
        $\True$ & $\True$ & $\False$ \\
        $\False$ & $\False$ & $\False$
      \end{tabular}
      \quad
      \begin{tabular}{c|cc}
        $\tf{\lor}$ & $\True$ & $\False$ \\
        \hline
        $\True$ & $\True$ & $\True$ \\
        $\False$ & $\True$ & $\False$
      \end{tabular}
      \quad
      \begin{tabular}{c|cc}
        $\tf{\lif}$ & $\True$ & $\False$ \\
        \hline
        $\True$ & $\True$ & $\False$ \\
        $\False$ & $\True$ & $\True$
      \end{tabular}
    \end{center}
    \caption{సాంప్రదాయిక తర్కం~$\LogCL$కు సత్యమూల్య ప్రమేయాలు.}
    \ollabel{fig:tf-CL}
  \end{figure}
  \end{ex}

\end{document}
